% =====================================================================
%  APÊNDICES -- Manual, seção 2.3.3
%  APÊNDICE = material ELABORADO PELO AUTOR (questionários, personas,
%  prompts, código-fonte). Pelo Manual, no sumário aparece só "APÊNDICES".
% =====================================================================
\begin{apendicesenv}
\partapendices
\addtocontents{toc}{\protect\setcounter{tocdepth}{-1}}

% ---------------------------------------------------------------------
\chapter{Especificação dos casos de uso}
\label{ap:casos-uso}

Este apêndice reproduz a especificação textual dos 30 casos de uso do
RoundUp, organizados pelas categorias da base de requisitos
(\autoref{sec:modelagem-funcional}). Quando o ator é o professor, a
academia herda o caso de uso por generalização de ator.

{\footnotesize\SingleSpacing
\setlength{\columnsep}{0.8cm}
\begin{multicols}{2}
% Gerado de RoundUp_Especificacao_Casos_de_Uso.docx
\section*{Contas}

\noindent\textbf{UC01 -- Cadastrar Conta}
\begin{itemize}[leftmargin=*, itemsep=0pt, topsep=2pt]
  \item \textbf{Atores:} Aluno, Professor, Academia
  \item \textbf{Requisitos relacionados:} RF01
  \item \textbf{Pré-condições:} O usuário não possui conta cadastrada no sistema.
  \item \textbf{Fluxo principal:}
  \begin{enumerate}[label=\arabic*., itemsep=0pt, topsep=0pt]
    \item O usuário acessa a tela de cadastro.
    \item O usuário seleciona o perfil desejado (Aluno, Professor ou Academia).
    \item O sistema exibe os campos obrigatórios correspondentes ao perfil escolhido.
    \item O usuário preenche os dados e confirma o cadastro.
    \item O sistema cria a conta e associa o perfil escolhido.
  \end{enumerate}
  \item \textbf{Fluxos alternativos:} A1. Dados inválidos ou incompletos: o sistema exibe mensagem de erro e mantém os dados já preenchidos.
  \item \textbf{Pós-condições:} Conta criada e usuário autenticado no perfil escolhido.
\end{itemize}
\medskip

\noindent\textbf{UC02 -- Autenticar-se e Recuperar Senha}
\begin{itemize}[leftmargin=*, itemsep=0pt, topsep=2pt]
  \item \textbf{Atores:} Aluno, Professor, Academia
  \item \textbf{Requisitos relacionados:} RF02, RNF03
  \item \textbf{Pré-condições:} O usuário possui conta cadastrada.
  \item \textbf{Fluxo principal:}
  \begin{enumerate}[label=\arabic*., itemsep=0pt, topsep=0pt]
    \item O usuário informa e-mail e senha.
    \item O sistema delega a validação das credenciais ao serviço de autenticação do Firebase.
    \item O sistema autentica o usuário e o redireciona à tela inicial do seu perfil.
  \end{enumerate}
  \item \textbf{Fluxos alternativos:} A1. Senha esquecida: o usuário solicita recuperação por e-mail; o sistema envia link de redefinição via Firebase. A2. Credenciais inválidas: o sistema exibe mensagem de erro.
  \item \textbf{Pós-condições:} Sessão iniciada.
\end{itemize}
\medskip

\noindent\textbf{UC03 -- Gerenciar Perfil da Academia}
\begin{itemize}[leftmargin=*, itemsep=0pt, topsep=2pt]
  \item \textbf{Atores:} Academia
  \item \textbf{Requisitos relacionados:} RF04
  \item \textbf{Pré-condições:} Academia autenticada.
  \item \textbf{Fluxo principal:}
  \begin{enumerate}[label=\arabic*., itemsep=0pt, topsep=0pt]
    \item A Academia acessa a tela de perfil institucional.
    \item A Academia edita nome, endereço, modalidades oferecidas e identidade visual.
    \item O sistema salva as alterações.
  \end{enumerate}
  \item \textbf{Fluxos alternativos:} A1. Dados inválidos: o sistema exibe mensagem de erro.
  \item \textbf{Pós-condições:} Perfil institucional atualizado.
\end{itemize}
\medskip

\noindent\textbf{UC04 -- Gerenciar Perfil do Professor}
\begin{itemize}[leftmargin=*, itemsep=0pt, topsep=2pt]
  \item \textbf{Atores:} Professor
  \item \textbf{Requisitos relacionados:} RF05
  \item \textbf{Pré-condições:} Professor autenticado.
  \item \textbf{Fluxo principal:}
  \begin{enumerate}[label=\arabic*., itemsep=0pt, topsep=0pt]
    \item O Professor acessa seu perfil.
    \item Edita nome, modalidades que leciona, graduação e indicação de atuação autônoma ou vinculada a academia.
    \item O sistema salva as alterações.
  \end{enumerate}
  \item \textbf{Fluxos alternativos:} A1. Dados inválidos: o sistema exibe mensagem de erro.
  \item \textbf{Pós-condições:} Perfil do professor atualizado.
\end{itemize}
\medskip

\noindent\textbf{UC05 -- Gerenciar Perfil do Aluno}
\begin{itemize}[leftmargin=*, itemsep=0pt, topsep=2pt]
  \item \textbf{Atores:} Aluno
  \item \textbf{Requisitos relacionados:} RF06
  \item \textbf{Pré-condições:} Aluno autenticado.
  \item \textbf{Fluxo principal:}
  \begin{enumerate}[label=\arabic*., itemsep=0pt, topsep=0pt]
    \item O Aluno acessa seu perfil.
    \item Visualiza e edita nome, modalidade praticada, faixa/grau atual e foto.
    \item O sistema salva as alterações.
  \end{enumerate}
  \item \textbf{Fluxos alternativos:} A1. Dados inválidos: o sistema exibe mensagem de erro.
  \item \textbf{Pós-condições:} Perfil do aluno atualizado.
\end{itemize}
\medskip

\section*{Vínculos}

\noindent\textbf{UC06 -- Solicitar Vínculo com Academia}
\begin{itemize}[leftmargin=*, itemsep=0pt, topsep=2pt]
  \item \textbf{Atores:} Professor
  \item \textbf{Requisitos relacionados:} RF07 , RN06
  \item \textbf{Pré-condições:} Professor autenticado.
  \item \textbf{Fluxo principal:}
  \begin{enumerate}[label=\arabic*., itemsep=0pt, topsep=0pt]
    \item O Professor busca e seleciona a academia desejada.
    \item O Professor envia a solicitação de vínculo.
    \item O sistema registra a solicitação com status Pendente e notifica a academia.
  \end{enumerate}
  \item \textbf{Fluxos alternativos:} A1. O Professor cancela a solicitação antes da resposta da academia.
  \item \textbf{Pós-condições:} Solicitação registrada, aguardando aceite da academia.
\end{itemize}
\medskip

\noindent\textbf{UC07 -- Gerenciar Vínculos de Professores}
\begin{itemize}[leftmargin=*, itemsep=0pt, topsep=2pt]
  \item \textbf{Atores:} Academia
  \item \textbf{Requisitos relacionados:} RF09 , RN06
  \item \textbf{Pré-condições:} Existem solicitações pendentes ou vínculos ativos com professores.
  \item \textbf{Fluxo principal:}
  \begin{enumerate}[label=\arabic*., itemsep=0pt, topsep=0pt]
    \item A Academia consulta a lista de solicitações de vínculo pendentes e vínculos ativos.
    \item A Academia aceita ou recusa uma solicitação, ou encerra um vínculo ativo.
    \item Ao aceitar, a Academia atribui as turmas sob responsabilidade do professor.
    \item O sistema atualiza o status do vínculo e notifica o professor.
  \end{enumerate}
  \item \textbf{Fluxos alternativos:} A1. Academia recusa a solicitação: sistema notifica o professor e encerra o fluxo.
  \item \textbf{Pós-condições:} Vínculo atualizado (ativo, recusado ou encerrado).
\end{itemize}
\medskip

\noindent\textbf{UC08 -- Vincular-se a Academia ou Professor}
\begin{itemize}[leftmargin=*, itemsep=0pt, topsep=2pt]
  \item \textbf{Atores:} Aluno
  \item \textbf{Requisitos relacionados:} RF08 , RN05
  \item \textbf{Pré-condições:} Aluno autenticado.
  \item \textbf{Fluxo principal:}
  \begin{enumerate}[label=\arabic*., itemsep=0pt, topsep=0pt]
    \item O Aluno busca uma academia ou um professor autônomo.
    \item O Aluno solicita o vínculo.
    \item O sistema registra o vínculo e define o escopo de turmas, missões, classificação e relatórios aplicáveis a esse aluno.
  \end{enumerate}
  \item \textbf{Fluxos alternativos:} A1. Aluno sem nenhum vínculo ativo: sistema bloqueia o acesso a turmas, sessões, missões e classificação (RN05).
  \item \textbf{Pós-condições:} Aluno vinculado, com escopo de uso definido.
\end{itemize}
\medskip

\section*{Treinos}

\noindent\textbf{UC09 -- Gerenciar Turmas}
\begin{itemize}[leftmargin=*, itemsep=0pt, topsep=2pt]
  \item \textbf{Atores:} Professor (Academia herda por generalização de ator)
  \item \textbf{Requisitos relacionados:} RF10 , RN07
  \item \textbf{Pré-condições:} Responsável autenticado (professor autônomo ou academia).
  \item \textbf{Fluxo principal:}
  \begin{enumerate}[label=\arabic*., itemsep=0pt, topsep=0pt]
    \item O responsável cria uma turma informando nome, modalidade, nível e professor responsável.
    \item O responsável edita ou encerra turmas existentes.
  \end{enumerate}
  \item \textbf{Fluxos alternativos:} A1. Tentativa de definir mais de um responsável para a mesma turma: sistema impede, pois toda turma tem responsável único (RN07).
  \item \textbf{Pós-condições:} Turma criada, atualizada ou encerrada.
\end{itemize}
\medskip

\noindent\textbf{UC10 -- Matricular-se em Turma}
\begin{itemize}[leftmargin=*, itemsep=0pt, topsep=2pt]
  \item \textbf{Atores:} Aluno
  \item \textbf{Requisitos relacionados:} RF11
  \item \textbf{Pré-condições:} Aluno possui vínculo ativo que dá acesso à turma.
  \item \textbf{Fluxo principal:}
  \begin{enumerate}[label=\arabic*., itemsep=0pt, topsep=0pt]
    \item O Aluno consulta as turmas disponíveis no escopo do seu vínculo.
    \item O Aluno se matricula em uma ou mais turmas.
  \end{enumerate}
  \item \textbf{Fluxos alternativos:} A1. Turma fora do escopo do vínculo do aluno: sistema não a exibe como disponível.
  \item \textbf{Pós-condições:} Matrícula registrada.
\end{itemize}
\medskip

\noindent\textbf{UC11 -- Criar Sessão de Treino}
\begin{itemize}[leftmargin=*, itemsep=0pt, topsep=2pt]
  \item \textbf{Atores:} Professor (Academia herda por generalização de ator)
  \item \textbf{Requisitos relacionados:} RF12
  \item \textbf{Pré-condições:} Existe turma cadastrada sob responsabilidade do solicitante.
  \item \textbf{Fluxo principal:}
  \begin{enumerate}[label=\arabic*., itemsep=0pt, topsep=0pt]
    \item O responsável seleciona a turma.
    \item Informa data, horário de início e término e o professor responsável pela sessão.
    \item Opcionalmente, gera sessões recorrentes a partir de uma grade semanal.
  \end{enumerate}
  \item \textbf{Fluxos alternativos:} A1. Conflito de horário com outra sessão da mesma turma: sistema alerta o responsável.
  \item \textbf{Pós-condições:} Sessão(ões) de treino criada(s).
\end{itemize}
\medskip

\noindent\textbf{UC12 -- Realizar Check-in}
\par\noindent Especificação apresentada no \autoref{qua:uc12-uc15}, de acordo com o comportamento implementado.

\medskip

\noindent\textbf{UC13 -- Consultar Agenda de Treinos}
\begin{itemize}[leftmargin=*, itemsep=0pt, topsep=2pt]
  \item \textbf{Atores:} Aluno
  \item \textbf{Requisitos relacionados:} RF14
  \item \textbf{Pré-condições:} Aluno matriculado em ao menos uma turma.
  \item \textbf{Fluxo principal:}
  \begin{enumerate}[label=\arabic*., itemsep=0pt, topsep=0pt]
    \item O Aluno acessa a agenda.
    \item O sistema exibe as próximas sessões das turmas em que está matriculado, com data, horário e professor responsável.
  \end{enumerate}
  \item \textbf{Pós-condições:} Agenda exibida.
\end{itemize}
\medskip

\noindent\textbf{UC14 -- Consultar Histórico de Presenças}
\begin{itemize}[leftmargin=*, itemsep=0pt, topsep=2pt]
  \item \textbf{Atores:} Aluno
  \item \textbf{Requisitos relacionados:} RF15
  \item \textbf{Pré-condições:} Aluno possui ao menos um \textit{check-in} validado.
  \item \textbf{Fluxo principal:}
  \begin{enumerate}[label=\arabic*., itemsep=0pt, topsep=0pt]
    \item O Aluno acessa o histórico de presenças.
    \item O sistema exibe os \textit{check-ins} validados ao longo do tempo.
  \end{enumerate}
  \item \textbf{Pós-condições:} Histórico exibido.
\end{itemize}
\medskip

\noindent\textbf{UC15 -- Validar Check-in}
\par\noindent Especificação apresentada no \autoref{qua:uc12-uc15}, de acordo com o comportamento implementado.

\medskip

\section*{Gamificação}

\noindent\textbf{UC16 -- Visualizar Nível e Progresso}
\begin{itemize}[leftmargin=*, itemsep=0pt, topsep=2pt]
  \item \textbf{Atores:} Aluno
  \item \textbf{Requisitos relacionados:} RF16 , RF17
  \item \textbf{Pré-condições:} Aluno autenticado.
  \item \textbf{Fluxo principal:}
  \begin{enumerate}[label=\arabic*., itemsep=0pt, topsep=0pt]
    \item O Aluno acessa seu painel.
    \item O sistema exibe o nível atual, o total de pontos acumulados (creditados a partir de \textit{check-ins} validados e missões concluídas) e quanto falta para o próximo nível.
  \end{enumerate}
  \item \textbf{Pós-condições:} Nível e progresso exibidos.
\end{itemize}
\medskip

\noindent\textbf{UC17 -- Consultar Missões}
\begin{itemize}[leftmargin=*, itemsep=0pt, topsep=2pt]
  \item \textbf{Atores:} Aluno
  \item \textbf{Requisitos relacionados:} RF18
  \item \textbf{Pré-condições:} Aluno vinculado a uma academia ou professor com missões ativas no seu escopo.
  \item \textbf{Fluxo principal:}
  \begin{enumerate}[label=\arabic*., itemsep=0pt, topsep=0pt]
    \item O Aluno acessa a lista de missões ativas no escopo do seu vínculo.
    \item O sistema exibe o progresso parcial de cada missão e sinaliza conclusão com a respectiva recompensa.
  \end{enumerate}
  \item \textbf{Pós-condições:} Missões e progresso exibidos.
\end{itemize}
\medskip

\noindent\textbf{UC18 -- Consultar Classificação}
\begin{itemize}[leftmargin=*, itemsep=0pt, topsep=2pt]
  \item \textbf{Atores:} Aluno
  \item \textbf{Requisitos relacionados:} RF20 , RN04 , RN09 , RNF11
  \item \textbf{Pré-condições:} Aluno optou por participar da classificação (adesão é opcional, conforme RNF11).
  \item \textbf{Fluxo principal:}
  \begin{enumerate}[label=\arabic*., itemsep=0pt, topsep=0pt]
    \item O Aluno acessa a tela de classificação.
    \item O sistema exibe o ranking por pontos obtidos na janela de tempo vigente (RN04), respeitando o escopo do vínculo do aluno -- academia, turma ou conjunto de alunos de um professor autônomo (RN09).
  \end{enumerate}
  \item \textbf{Fluxos alternativos:} A1. Aluno não participante: a tela informa que a classificação está desativada e oferece a opção de habilitar a adesão.
  \item \textbf{Pós-condições:} Classificação exibida, ou aviso de não adesão.
\end{itemize}
\medskip

\noindent\textbf{UC19 -- Consultar Sequência de Treinos}
\begin{itemize}[leftmargin=*, itemsep=0pt, topsep=2pt]
  \item \textbf{Atores:} Aluno
  \item \textbf{Requisitos relacionados:} RF21
  \item \textbf{Pré-condições:} Aluno autenticado.
  \item \textbf{Fluxo principal:}
  \begin{enumerate}[label=\arabic*., itemsep=0pt, topsep=0pt]
    \item O Aluno acessa seu painel.
    \item O sistema exibe a sequência de treinos consecutivos do aluno.
  \end{enumerate}
  \item \textbf{Pós-condições:} Sequência exibida.
\end{itemize}
\medskip

\section*{Gestão}

\noindent\textbf{UC20 -- Criar Missão Personalizada}
\begin{itemize}[leftmargin=*, itemsep=0pt, topsep=2pt]
  \item \textbf{Atores:} Professor (Academia herda por generalização de ator)
  \item \textbf{Requisitos relacionados:} RF23
  \item \textbf{Pré-condições:} Responsável autenticado.
  \item \textbf{Fluxo principal:}
  \begin{enumerate}[label=\arabic*., itemsep=0pt, topsep=0pt]
    \item O responsável acessa a criação de missões.
    \item Define título, critério de conclusão, recompensa em pontos, público-alvo e período de vigência.
    \item O sistema publica a missão e notifica os alunos do público-alvo (RF31).
  \end{enumerate}
  \item \textbf{Fluxos alternativos:} A1. Período de vigência inválido: sistema exibe erro e não publica a missão.
  \item \textbf{Pós-condições:} Missão publicada e visível aos alunos do público-alvo (UC17).
\end{itemize}
\medskip

\noindent\textbf{UC21 -- Consultar Alunos e Frequência}
\begin{itemize}[leftmargin=*, itemsep=0pt, topsep=2pt]
  \item \textbf{Atores:} Professor (Academia herda por generalização de ator)
  \item \textbf{Requisitos relacionados:} RF24
  \item \textbf{Pré-condições:} Responsável possui ao menos um aluno vinculado.
  \item \textbf{Fluxo principal:}
  \begin{enumerate}[label=\arabic*., itemsep=0pt, topsep=0pt]
    \item O responsável acessa a lista de alunos sob sua responsabilidade.
    \item O sistema exibe a frequência recente de cada aluno.
  \end{enumerate}
  \item \textbf{Pós-condições:} Lista e frequência exibidas.
\end{itemize}
\medskip

\noindent\textbf{UC22 -- Registrar Graduação de Faixa}
\begin{itemize}[leftmargin=*, itemsep=0pt, topsep=2pt]
  \item \textbf{Atores:} Professor (Academia herda por generalização de ator)
  \item \textbf{Requisitos relacionados:} RF25 , RN03
  \item \textbf{Pré-condições:} Aluno vinculado ao responsável.
  \item \textbf{Fluxo principal:}
  \begin{enumerate}[label=\arabic*., itemsep=0pt, topsep=0pt]
    \item O responsável seleciona o aluno.
    \item Registra a nova faixa/grau concedida.
    \item O sistema atualiza o perfil do aluno e o histórico de progressão.
  \end{enumerate}
  \item \textbf{Pós-condições:} Graduação registrada. A concessão é sempre uma decisão humana do responsável, nunca automática por acúmulo de pontos (RN03).
\end{itemize}
\medskip

\noindent\textbf{UC23 -- Gerar Relatório de Frequência}
\begin{itemize}[leftmargin=*, itemsep=0pt, topsep=2pt]
  \item \textbf{Atores:} Professor (Academia herda por generalização de ator)
  \item \textbf{Requisitos relacionados:} RF26
  \item \textbf{Pré-condições:} Existem sessões e \textit{check-ins} registrados no período selecionado.
  \item \textbf{Fluxo principal:}
  \begin{enumerate}[label=\arabic*., itemsep=0pt, topsep=0pt]
    \item O responsável seleciona a turma e o período desejado.
    \item O sistema gera o relatório de frequência.
  \end{enumerate}
  \item \textbf{Pós-condições:} Relatório gerado, disponível também como fonte de dados para a avaliação do TCC.
\end{itemize}
\medskip

\noindent\textbf{UC24 -- Consultar Visão Consolidada}
\begin{itemize}[leftmargin=*, itemsep=0pt, topsep=2pt]
  \item \textbf{Atores:} Academia (uso exclusivo)
  \item \textbf{Requisitos relacionados:} RF27
  \item \textbf{Pré-condições:} Academia autenticada, com turmas, professores e alunos cadastrados.
  \item \textbf{Fluxo principal:}
  \begin{enumerate}[label=\arabic*., itemsep=0pt, topsep=0pt]
    \item A Academia acessa o painel consolidado.
    \item O sistema exibe todas as turmas, professores e alunos da academia, com indicadores agregados de frequência.
  \end{enumerate}
  \item \textbf{Pós-condições:} Visão consolidada exibida.
\end{itemize}
\medskip

\noindent\textbf{UC25 -- Personalizar Pontuação e Graduação}
\begin{itemize}[leftmargin=*, itemsep=0pt, topsep=2pt]
  \item \textbf{Atores:} Professor (Academia herda por generalização de ator)
  \item \textbf{Requisitos relacionados:} RF33
  \item \textbf{Pré-condições:} Responsável autenticado.
  \item \textbf{Fluxo principal:}
  \begin{enumerate}[label=\arabic*., itemsep=0pt, topsep=0pt]
    \item O responsável acessa as configurações de gamificação do seu perfil.
    \item Personaliza a quantidade de pontos por presença, os critérios de subida de nível e a estrutura de faixas da modalidade.
    \item O sistema salva a configuração (ConfiguracaoGamificacao) associada ao perfil.
  \end{enumerate}
  \item \textbf{Fluxos alternativos:} A1. Responsável não personaliza: o sistema aplica a configuração padrão pré-definida.
  \item \textbf{Pós-condições:} Configuração de gamificação salva e aplicada aos alunos do escopo do responsável.
\end{itemize}
\medskip

\section*{Conformidade}

\noindent\textbf{UC26 -- Registrar Consentimento LGPD}
\begin{itemize}[leftmargin=*, itemsep=0pt, topsep=2pt]
  \item \textbf{Atores:} Aluno, Professor, Academia
  \item \textbf{Requisitos relacionados:} RF28 , RNF10
  \item \textbf{Pré-condições:} Usuário está realizando o cadastro (UC01).
  \item \textbf{Fluxo principal:}
  \begin{enumerate}[label=\arabic*., itemsep=0pt, topsep=0pt]
    \item O sistema apresenta a política de privacidade durante o cadastro.
    \item O usuário registra o consentimento quanto ao tratamento de dados pessoais.
  \end{enumerate}
  \item \textbf{Fluxos alternativos:} A1. Usuário não consente: cadastro não é concluído.
  \item \textbf{Pós-condições:} Consentimento registrado, associado à conta.
\end{itemize}
\medskip

\noindent\textbf{UC27 -- Solicitar Exclusão de Conta e Dados}
\begin{itemize}[leftmargin=*, itemsep=0pt, topsep=2pt]
  \item \textbf{Atores:} Aluno, Professor, Academia
  \item \textbf{Requisitos relacionados:} RF29
  \item \textbf{Pré-condições:} Usuário autenticado.
  \item \textbf{Fluxo principal:}
  \begin{enumerate}[label=\arabic*., itemsep=0pt, topsep=0pt]
    \item O usuário acessa as configurações da conta.
    \item Solicita a exclusão da conta e dos dados pessoais associados.
    \item O sistema processa a exclusão conforme a LGPD (Lei 13.709/2018).
  \end{enumerate}
  \item \textbf{Pós-condições:} Conta e dados pessoais eliminados, conforme direito de eliminação previsto na LGPD.
\end{itemize}
\medskip

\noindent\textbf{UC28 -- Registrar Responsável Legal}
\begin{itemize}[leftmargin=*, itemsep=0pt, topsep=2pt]
  \item \textbf{Atores:} Aluno menor de 18 anos
  \item \textbf{Requisitos relacionados:} RF30
  \item \textbf{Pré-condições:} Aluno identificado como menor de idade no cadastro.
  \item \textbf{Fluxo principal:}
  \begin{enumerate}[label=\arabic*., itemsep=0pt, topsep=0pt]
    \item O sistema solicita a identificação do responsável legal.
    \item O responsável legal fornece o consentimento necessário.
    \item O sistema registra o vínculo de responsabilidade.
  \end{enumerate}
  \item \textbf{Pós-condições:} Responsável legal registrado e vinculado à conta do aluno menor.
\end{itemize}
\medskip

\section*{Retenção}

\noindent\textbf{UC29 -- Receber Notificações}
\begin{itemize}[leftmargin=*, itemsep=0pt, topsep=2pt]
  \item \textbf{Atores:} Aluno
  \item \textbf{Requisitos relacionados:} RF31
  \item \textbf{Pré-condições:} Aluno com conta ativa e notificações habilitadas.
  \item \textbf{Fluxo principal:}
  \begin{enumerate}[label=\arabic*., itemsep=0pt, topsep=0pt]
    \item O sistema envia notificação de lembrete de sessão de treino, publicação de nova missão ou conquista obtida.
    \item O Aluno visualiza a notificação no aplicativo.
  \end{enumerate}
  \item \textbf{Pós-condições:} Notificação entregue e marcada como lida ao ser aberta.
\end{itemize}
\medskip

\noindent\textbf{UC30 -- Consultar Alertas de Queda de Frequência}
\begin{itemize}[leftmargin=*, itemsep=0pt, topsep=2pt]
  \item \textbf{Atores:} Professor (Academia herda por generalização de ator)
  \item \textbf{Requisitos relacionados:} RF32
  \item \textbf{Pré-condições:} Existe aluno sob responsabilidade do solicitante com queda relevante de frequência no período.
  \item \textbf{Fluxo principal:}
  \begin{enumerate}[label=\arabic*., itemsep=0pt, topsep=0pt]
    \item O sistema sinaliza, no painel do responsável, os alunos com queda relevante de frequência.
    \item O responsável consulta o alerta e pode tomar uma ação, como contatar o aluno.
  \end{enumerate}
  \item \textbf{Pós-condições:} Alerta consultado. É o requisito que mais dialoga diretamente com a hipótese de redução da evasão.
\end{itemize}
\medskip


\end{multicols}
}

% ---------------------------------------------------------------------
\chapter{Questionário de avaliação das ideias}
\label{ap:questionario}

Este apêndice reproduz o questionário aplicado aos usuários reais, também
usado na geração das respostas sintéticas. O formulário foi feito no
Google Forms, com tempo estimado de resposta de 6 a 8 minutos, e abria com
o texto de apresentação e consentimento abaixo.

\begin{citacao}
Este questionário faz parte de um Trabalho de Conclusão de Curso em
Engenharia da Computação no Centro Universitário Ibmec. O objetivo é
entender o que motiva as pessoas a continuarem treinando lutas e artes
marciais, o que leva ao abandono da prática, e avaliar a receptividade a
um aplicativo de acompanhamento e engajamento voltado a academias.

As respostas são anônimas e serão utilizadas exclusivamente para fins
acadêmicos, de forma agregada, sem identificação individual. O contato
(quando fornecido) é guardado separadamente das respostas e usado só para
convite a testes. O fornecimento de contato ao final é opcional e serve
apenas para convite a testes futuros do aplicativo. Você pode desistir a
qualquer momento simplesmente fechando a página. O tratamento dos dados
segue a Lei Geral de Proteção de Dados (Lei 13.709/2018).

Em caso de dúvidas, entre em contato pelo e-mail
andreccoelho2004@gmail.com.
\end{citacao}

Depois das perguntas de perfil, o respondente era encaminhado a uma das
três seções conforme a resposta à quarta pergunta, e todos terminavam na
seção de encerramento. As alternativas das perguntas fechadas foram
reconstruídas a partir das respostas registradas no formulário e das
listas do \textit{script} de geração das respostas sintéticas, que
validava cada resposta contra as alternativas do formulário; por isso,
todas as alternativas listadas existem no formulário, mas alguma
alternativa nunca escolhida pode não aparecer.

{\small\SingleSpacing
\setlength{\columnsep}{0.8cm}
\begin{multicols}{2}
% Gerado a partir do cabecalho das planilhas de respostas, com as alternativas observadas.
\section*{Termo de consentimento e perfil (todos os respondentes)}
\begin{enumerate}[leftmargin=*, itemsep=1pt, parsep=0pt]
  \item Li as informações acima e concordo em participar da pesquisa.\\
        {\footnotesize\textit{Escolha única:} Sim, concordo em participar.}
  \item Faixa etária\\
        {\footnotesize\textit{Escolha única:} Menos que 18 anos; 18 a 24 anos; 25 a 34 anos; 35 a 44 anos; 45 a 54 anos; 55 anos ou mais.}
  \item Gênero\\
        {\footnotesize\textit{Escolha única:} Feminino; Masculino; Prefiro não dizer.}
  \item Qual das opções descreve melhor a sua relação com lutas ou artes marciais hoje?\\
        {\footnotesize\textit{Escolha única:} Treino atualmente; Já treinei, mas parei; Sou professor ou instrutor.}
\end{enumerate}

\section*{Seção do professor}
\begin{enumerate}[leftmargin=*, itemsep=1pt, parsep=0pt]
  \item Há quanto tempo você dá aulas de luta ou arte marcial?\\
        {\footnotesize\textit{Escolha única:} De 1 a 3 anos; De 3 a 5 anos; De 5 a 10 anos; Mais de 10 anos.}
  \item Quais modalidades você ensina?\\
        {\footnotesize\textit{Múltipla escolha:} Boxe; Muay Thai; Jiu-Jitsu; Karatê; Taekwondo; Judô; MMA.}
  \item Aproximadamente quantos alunos ativos você tem hoje? (na academia e particulares)\\
        {\footnotesize\textit{Escolha única:} Até 20; De 21 a 50; De 51 a 100; De 101 a 200; Mais de 200.}
  \item De cada 10 alunos que se matriculam, quantos ainda estão treinando depois de 6 meses?\\
        {\footnotesize\textit{Resposta numérica de 0 a 10.}}
  \item Em que momento a desistência costuma ser mais frequente?\\
        {\footnotesize\textit{Escolha única:} No primeiro mês; Entre 1 e 3 meses; Entre 3 e 6 meses; Não há um padrão claro.}
  \item Na sua percepção, quais são os principais motivos que levam os alunos a abandonarem o treino? (Selecione até 3 mais relevantes)\\
        {\footnotesize\textit{Múltipla escolha:} Falta de tempo ou conflito com a rotina; Questões financeiras; Não perceber evolução técnica; Ausência de metas claras; Lesão ou problema de saúde; Vergonha, insegurança ou intimidação; Falta de vínculo com a turma; Horários das aulas; Cansaço e desgaste da rotina de treino.}
  \item Como você acompanha a frequência dos seus alunos atualmente?\\
        {\footnotesize\textit{Escolha única:} Aplicativo de mensagens (WhatsApp); Caderno ou lista em papel; Sistema de gestão da academia; Planilha (Excel, Google Sheets); Não acompanho de forma sistemática.}
  \item Como funciona a graduação ou troca de faixa na sua academia, e como você controla quem está apto a graduar?\\
        {\footnotesize\textit{Resposta aberta.}}
  \item Quando um aluno some por algumas semanas, você costuma entrar em contato?\\
        {\footnotesize\textit{Escolha única:} Sim, sempre; Às vezes, quando percebo; Raramente.}
  \item Qual é hoje a sua maior dificuldade na gestão da turma?\\
        {\footnotesize\textit{Resposta aberta.}}
  \item Avalie de 1 (pouco útil) a 5 (muito útil) cada recurso:
  \begin{enumerate}[label=\alph*)]
    \item Pontos de experiência (XP) atribuídos automaticamente a cada treino
    \item Registro rápido de presença (check-in) a cada treino
    \item Contador de dias treinados e frequência ao longo das semanas
    \item Missões e desafios semanais criados por você para a turma
    \item Ranking de frequência entre os alunos da turma
    \item Medalhas e conquistas por marcos atingidos (100 treinos, 6 meses seguidos)
    \item Barra de progresso rumo à próxima graduação, visível para o aluno
    \item Alerta automático quando um aluno fica mais de 15 dias sem aparecer
    \item Relatório de frequência e evolução da turma
  \end{enumerate}
  \item Que riscos ou problemas você vê no uso de ranking e competição entre alunos dentro da academia?\\
        {\footnotesize\textit{Resposta aberta.}}
  \item Quanto tempo por semana você estaria disposto a dedicar para alimentar um aplicativo desse tipo?\\
        {\footnotesize\textit{Escolha única:} Menos de 5 minutos; De 5 a 15 minutos; De 15 a 30 minutos; Mais de 30 minutos.}
  \item Você utiliza algum sistema ou aplicativo de gestão hoje? Se sim, qual, e o que falta nele?\\
        {\footnotesize\textit{Resposta aberta.}}
  \item Se você pudesse pedir uma única funcionalidade para esse aplicativo, qual seria?\\
        {\footnotesize\textit{Resposta aberta.}}
  \item O que faria você desistir de usar o aplicativo depois de algumas semanas?\\
        {\footnotesize\textit{Resposta aberta.}}
  \item Existe alguma prática da sua academia que hoje funciona bem para segurar aluno e que poderia virar funcionalidade no aplicativo?\\
        {\footnotesize\textit{Resposta aberta.}}
\end{enumerate}

\section*{Seção do praticante}
\begin{enumerate}[leftmargin=*, itemsep=1pt, parsep=0pt]
  \item Qual modalidade você pratica atualmente?\\
        {\footnotesize\textit{Múltipla escolha:} Boxe; Muay Thai; Jiu-Jitsu; Karatê; Taekwondo; Judô; MMA.}
  \item Há quanto tempo você treina?\\
        {\footnotesize\textit{Escolha única:} Menos de 3 meses; De 3 a 6 meses; De 6 meses a 1 ano; De 1 a 3 anos; Mais de 3 anos.}
  \item Quantas vezes por semana você treina, em média?\\
        {\footnotesize\textit{Escolha única:} 1 vez; 2 vezes; 3 vezes; 4 vezes; 5 vezes ou mais.}
  \item O que faz você continuar treinando? (Selecione até 3)\\
        {\footnotesize\textit{Múltipla escolha:} Saúde física e condicionamento; Saúde mental e alívio do estresse; Amizades e convívio com a turma; Defesa pessoal; Disciplina e rotina; Competição; Evolução técnica e aprendizado; Estética e composição corporal; Relação com o professor.}
  \item De 1 a 5, o quanto você consegue perceber a sua própria evolução técnica?\\
        {\footnotesize\textit{Escala de 1 a 5.}}
  \item Você sabe exatamente o que precisa fazer para chegar à próxima graduação ou faixa?\\
        {\footnotesize\textit{Escolha única:} Sim, sei com clareza; Não faço ideia; Tenho uma ideia aproximada; Minha modalidade ou academia não trabalha com graduação.}
  \item Você já pensou em parar de treinar em algum momento? O que quase o/a fez você parar?\\
        {\footnotesize\textit{Resposta aberta.}}
  \item Você já parou de treinar e voltou depois?\\
        {\footnotesize\textit{Escolha única:} Sim; Não.}
  \item Você acompanha a sua frequência de treinos de alguma forma?\\
        {\footnotesize\textit{Escolha única:} Só mentalmente, sem registro; Não acompanho; Sim, em anotações próprias ou planilha; Sim, em aplicativo.}
  \item Avalie de 1 (pouco útil) a 5 (muito útil) cada recurso:
  \begin{enumerate}[label=\alph*)]
    \item Pontos de experiência (XP) ganhos a cada treino
    \item Registro rápido de presença (check-in) a cada treino
    \item Contador de dias treinados e frequência ao longo das semanas
    \item Sequência de dias ou semanas treinadas sem quebrar (ofensiva)
    \item Missões e desafios semanais propostos pelo professor
    \item Ranking de frequência entre os alunos da turma
    \item Medalhas por marcos atingidos (100 treinos, 1 ano de prática)
    \item Barra de progresso mostrando o quanto falta para a próxima graduação
    \item Histórico completo de presença e evolução pessoal
    \item Mural ou feed da turma, com conquistas dos colegas
    \item Lembretes automáticos nos dias em que você costuma treinar
  \end{enumerate}
  \item De 1 a 5, o quanto você usaria um aplicativo assim se os seus dados fossem visíveis apenas para você e para o seu professor?\\
        {\footnotesize\textit{Escala de 1 a 5.}}
  \item Você já usou algum aplicativo com pontos, metas ou sequências que ajudou a manter um hábito? Qual, e o que funcionou nele?\\
        {\footnotesize\textit{Resposta aberta.}}
  \item Se esse aplicativo existisse, que funcionalidade você gostaria que ele tivesse?\\
        {\footnotesize\textit{Resposta aberta.}}
  \item O que faria você desinstalar o aplicativo?\\
        {\footnotesize\textit{Resposta aberta.}}
  \item Que informação sobre o seu treino você gostaria de ver registrada e que hoje se perde?\\
        {\footnotesize\textit{Resposta aberta.}}
\end{enumerate}

\section*{Seção do ex-praticante}
\begin{enumerate}[leftmargin=*, itemsep=1pt, parsep=0pt]
  \item Qual modalidade você praticava?\\
        {\footnotesize\textit{Escolha única:} Boxe; Muay Thai; Jiu-Jitsu; Karatê; Taekwondo; Judô; MMA.}
  \item Por quanto tempo você treinou antes de parar?\\
        {\footnotesize\textit{Escolha única:} Menos de 1 mês; De 1 a 3 meses; De 3 a 6 meses; De 6 meses a 1 ano; De 1 a 3 anos; Mais de 3 anos.}
  \item Há quanto tempo você parou?\\
        {\footnotesize\textit{Escolha única:} Menos de 6 meses; De 6 meses a 1 ano; De 1 a 3 anos; Mais de 3 anos.}
  \item Quantas vezes por semana você treinava, em média?\\
        {\footnotesize\textit{Escolha única:} 1 vez; 2 vezes; 3 vezes; 4 vezes; 5 vezes ou mais.}
  \item Quais fatores contribuíram para você parar de treinar? (Selecione até 3)\\
        {\footnotesize\textit{Múltipla escolha:} Falta de tempo ou conflito com a rotina; Questões financeiras; Não perceber evolução técnica; Ausência de metas claras; Lesão ou problema de saúde; Vergonha, insegurança ou intimidação; Falta de vínculo com a turma; Pouca proximidade com o professor; Horários das aulas; Mudança de cidade ou bairro; Cansaço e desgaste da rotina de treino; Perda de interesse pela modalidade.}
  \item Entre os fatores que você marcou, qual foi o principal?\\
        {\footnotesize\textit{Escolha única:} Falta de tempo ou conflito com a rotina; Questões financeiras; Não perceber evolução técnica; Ausência de metas claras; Lesão ou problema de saúde; Vergonha, insegurança ou intimidação; Falta de vínculo com a turma; Pouca proximidade com o professor; Horários das aulas; Mudança de cidade ou bairro; Cansaço e desgaste da rotina de treino; Perda de interesse pela modalidade.}
  \item Conte com suas palavras como foi o processo até você parar de treinar\\
        {\footnotesize\textit{Resposta aberta.}}
  \item De 1 a 5, o quanto você sentia que estava evoluindo tecnicamente pouco antes de parar?\\
        {\footnotesize\textit{Escala de 1 a 5.}}
  \item Você sabia o que precisava fazer para chegar à próxima graduação?\\
        {\footnotesize\textit{Escolha única:} Sim, sabia com clareza; Tinha uma ideia aproximada; Não fazia ideia.}
  \item Depois que você parou de aparecer, alguém da academia tentou te chamar de volta?\\
        {\footnotesize\textit{Escolha única:} Sim, o professor ou instrutor me chamou de volta; Ninguém entrou em contato; Sim, a administração ou secretaria entrou em contato; Alguém da academia falou comigo, mas apenas socialmente, sem me chamar para voltar.}
  \item Se alguém entrou em contato, isso fez você considerar voltar?\\
        {\footnotesize\textit{Escolha única:} Sim, considerei, mas não voltei; Não fez diferença; Não se aplica, ninguém me chamou de volta; Sim, cheguei a voltar.}
  \item Você tem vontade de voltar a treinar?\\
        {\footnotesize\textit{Escolha única:} Sim, pretendo voltar em breve; Talvez, um dia; Não pretendo voltar.}
  \item O que faria você voltar a treinar hoje?\\
        {\footnotesize\textit{Resposta aberta.}}
  \item Avalie de 1 (pouco útil) a 5 (muito útil) cada recurso:
  \begin{enumerate}[label=\alph*)]
    \item Pontos de experiência (XP) ganhos a cada treino
    \item Registro rápido de presença (check-in) a cada treino
    \item Contador de dias treinados e frequência ao longo das semanas
    \item Sequência de dias ou semanas treinadas sem quebrar (ofensiva)
    \item Missões e desafios semanais propostos pelo professor
    \item Ranking de frequência entre os alunos da turma
    \item Medalhas por marcos atingidos
    \item Barra de progresso mostrando o quanto falta para a próxima graduação
    \item Histórico completo de presença e evolução pessoal
    \item Mural ou feed da turma, com conquistas dos colegas
    \item Mensagem automática do professor após alguns dias de ausência
    \item Lembretes automáticos nos dias de treino
  \end{enumerate}
  \item Sendo sincero, você acredita que um aplicativo assim teria mudado a sua decisão de parar?\\
        {\footnotesize\textit{Escolha única:} Não, o motivo pelo qual parei não tem relação com isso; Talvez, teria ajudado, mas não sei se seria suficiente; Sim, provavelmente eu teria continuado.}
  \item Algum desses recursos teria te incomodado ou te afastado?\\
        {\footnotesize\textit{Escolha única:} Nenhum me incomodaria; Lembretes e notificações; Mural com conquistas dos colegas; Ranking visível para outros alunos; Mensagem do professor cobrando presença.}
  \item Que funcionalidade um aplicativo da academia precisaria ter para fazer você voltar e permanecer?\\
        {\footnotesize\textit{Resposta aberta.}}
  \item O que faria você desinstalar um aplicativo desse tipo?\\
        {\footnotesize\textit{Resposta aberta.}}
\end{enumerate}

\section*{Encerramento (todos os respondentes)}
\begin{enumerate}[leftmargin=*, itemsep=1pt, parsep=0pt]
  \item Deseja acrescentar algum comentário sobre o tema?\\
        {\footnotesize\textit{Resposta aberta.}}
  \item Você aceitaria participar de um teste do aplicativo ou de uma conversa rápida de acompanhamento?\\
        {\footnotesize\textit{Escolha única:} Sim; Não.}
  \item Deixe um e-mail ou telefone para contato (opcional).\\
        {\footnotesize\textit{Resposta aberta.}}
\end{enumerate}


\end{multicols}
}

% ---------------------------------------------------------------------
\chapter{Personas sintéticas}
\label{ap:personas}

As 12 personas abaixo são sintéticas. Elas foram escritas a partir do
padrão das seis respostas reais e serviram para apoiar decisões de
experiência do usuário e para orientar as regras do gerador de respostas
de teste (\autoref{sec:val-personas}). As personas de professor e de dono
de academia são hipóteses a validar, porque nenhum professor real
respondeu ao questionário. O \autoref{qua:mapa-personas} resume as
personas.

\begin{quadro}[!htb]
  \caption{Mapa das personas sintéticas}
  \label{qua:mapa-personas}
  \centering
  \footnotesize
  \begin{tabular}{|c|p{2.3cm}|p{2.6cm}|c|p{1.9cm}|p{3.4cm}|}
    \hline
    \textbf{ID} & \textbf{Persona} & \textbf{Perfil} & \textbf{Faixa} & \textbf{Modalidade} & \textbf{Postura com a gamificação} \\ \hline
    A1 & Davi, 16 & Aluno ativo & $<$ 18 & Jiu-Jitsu & Gosta, mas responde no automático \\ \hline
    A2 & Lucas, 21 & Aluno ativo & 18--24 & Muay Thai & Entusiasta de \textit{ranking} e ofensiva \\ \hline
    A3 & Rafaela, 29 & Aluno ativo & 25--34 & Boxe & Quer dados, rejeita exposição \\ \hline
    A4 & Patrícia, 38 & Aluno ativo & 35--44 & Jiu-Jitsu & Aprova tudo de forma morna \\ \hline
    A5 & Marcelo, 42 & Aluno ativo & 35--44 & Jiu-Jitsu & Cético, vê o app como trabalho extra \\ \hline
    A6 & Sérgio, 51 & Aluno ativo & 45+ & Muay Thai & Indiferente, valoriza só o histórico \\ \hline
    E1 & Gabriel, 20 & Ex-aluno & 18--24 & Jiu-Jitsu & Gosta da ideia, odeia notificação \\ \hline
    E2 & Thiago, 31 & Ex-aluno & 25--34 & Boxe & Neutro, nota 4 em tudo \\ \hline
    E3 & Juliana, 40 & Ex-aluno & 35--44 & Muay Thai & Acha que lembrete teria ajudado \\ \hline
    E4 & Roberto, 48 & Ex-aluno & 45+ & Jiu-Jitsu & Resistente, prefere contato humano \\ \hline
    P1 & Felipe, 27 & Professor & 25--34 & Muay Thai & Quer alerta de ausência e relatório \\ \hline
    P2 & Cláudio, 49 & Professor e dono de academia & 45+ & Boxe e Muay Thai & Teme \textit{ranking}, quer controle de graduação \\ \hline
  \end{tabular}
  \fonte{Elaborado pelo autor (2026).}
\end{quadro}

As fichas completas, com rotina, motivações, riscos de evasão e
respostas-modelo de cada persona, estão no documento de personas do
projeto.

% ---------------------------------------------------------------------
\chapter{Protocolo de geração das respostas sintéticas}
\label{ap:prompt}

\begin{quadro}[!htb]
  \caption{Protocolo de geração e envio das respostas sintéticas}
  \label{qua:protocolo}
  \centering
  \footnotesize
  \begin{tabular}{|p{3.4cm}|p{9.8cm}|}
    \hline
    \textbf{Item} & \textbf{Descrição} \\ \hline
    Ferramenta de IA & Claude (Anthropic); identificador do modelo configurado na sessão: \texttt{claude-opus-5-5} \\ \hline
    Data da sessão & 26/10/2026 \\ \hline
    Objetivo & Testar o formulário e a análise sem misturar dados sintéticos com a coleta real \\ \hline
    Base de referência & Seis respostas reais coletadas em 18 e 19/09/2026, usadas para extrair o padrão de escrita \\ \hline
    Personas & 12 fichas por perfil (aluno ativo, ex-aluno, professor) e faixa etária (\autoref{ap:personas}) \\ \hline
    Método de geração & \textit{Script} em Python com bancos de respostas curtas por pergunta e regras por perfil e idade; combinações sorteadas pelo \textit{script}. Não houve um \textit{prompt} por resposta \\ \hline
    Parâmetros & Semente \texttt{random.seed(26102026)} e distribuição da \autoref{tab:perfil-personas}; sem temperatura ou outro parâmetro do modelo, pois o sorteio foi feito pelo \textit{script} \\ \hline
    Relação com as fichas & As 12 fichas orientaram as regras; as 65 linhas não correspondem uma a uma às fichas \\ \hline
    Instruções ao modelo & Dadas em conversa; o registro reproduzível são os \textit{scripts} versionados no repositório (ver abaixo) \\ \hline
    Ambiente de envio & Cópia do formulário (``RoundUp (Personas IA)''), sem planilha vinculada; nenhuma resposta sintética entrou na base da pesquisa \\ \hline
    Método de envio & Google Apps Script no editor da cópia: validação contra as alternativas do formulário e registro via FormApp (\texttt{createResponse}), sem passar pelo link público \\ \hline
    Janela de envio & 26/10/2026, de 13h12 a 13h14 \\ \hline
    Identificação & Campo de contato com \texttt{sintetico-<ID>}: A01 a A32, E01 a E23, P01 a P10 \\ \hline
    Total enviado & 65 respostas, todas com o marcador, conferidas na exportação \\ \hline
  \end{tabular}
  \fonte{Elaborado pelo autor (2026).}
\end{quadro}

Os \textit{scripts} de geração e de envio estão no repositório do projeto,
em \texttt{docs/pesquisa/respostas-sinteticas/}
(\url{https://github.com/andreccoelho/roundup_app/tree/dev/docs/pesquisa/respostas-sinteticas}):
\texttt{gerar\_respostas.py}, que gera as respostas com a semente
26102026, e \texttt{enviar\_respostas.gs}, que valida e envia as respostas
pelo Google Apps Script. O diretório traz também um \texttt{README.md} com
as instruções de reprodução.

\end{apendicesenv}
